\clearpage
\fancyhead[L]{\small VESTA.}
\hypertarget{native-295-lower}{}
\centerline{VESTA.}

\begin{center}
\rule{0.5\linewidth}{0.4pt}\\[2pt]
Monatliche Correspondenz. Band XVI. S. 285–291. 1807 September.\\[2pt]
\rule{0.5\linewidth}{0.4pt}
\end{center}

\smallskip
\noindent---\ ---\ ---

\smallskip
\par\indent „Mit dem 8. Juli, schreibt uns Dr. \textsc{Olbers}, habe ich meine Beobachtungen geschlossen, da der
niedere Stand des Planeten, Mondschein und Dämmerung sie zu sehr erschwerten. Die sechs letztern
Beobachtungen sind in Gegenwart unseres vortrefflichen Freundes, \textsc{Gauss}, angestellt worden, der mich
mit einem Besuch erfreut hat. Dieser vortreffliche Astronom wird uns nun sehr bald die III. Elemente
der \emph{Vesta} geben. Indessen werden die zweiten Elemente nur unbedeutende Correctionen erleiden, da sie
auch die letzten Beobachtungen noch so genau darstellen, und die AR. nur etwa $1'$ zu gross, die Declinat.
noch keine $30''$ zu klein geben. Nach den II. Elementen habe ich gefunden, dass \emph{Vesta} am 25. März
1796 in der Zone stand, die damals nach der \emph{Hist. Cél.} beobachtet wurde. Doch kann ein kleiner sehr
\SourceLeaf{296}möglicher Fehler in der halben grossen Achse, und also in der mittlern Bewegung der \emph{Vesta}, die Declination derselben für diesen Tag noch sehr ändern. Ob ein beobachteter Stern in dieser Zone jetzt fehlt,
weiss ich nicht, weil der Theil derselben, in dem die \emph{Vesta} stehen konnte, jetzt unter den Sonnenstrahlen verborgen ist.“

\par\indent Diese \emph{dritten Elemente}, zu denen Dr. \textsc{Gauss} ausser allen in dieser Zeitschrift schon bekannt gemachten Beobachtungen auch noch eine Reihe Pariser Beobachtungen vom 13. April bis 22. Mai, die ihm
Hr. \textsc{Bouvard} mitgetheilt, benutzt hatte, erhielten wir vor wenig Tagen, und wir lassen solche hier
folgen:

\smallskip

\begin{center}\begin{tabular}{lr}
Epoche 31. März 1807 Mittag in Bremen  & $192^\circ\;23'\;30''1$ \\
Tägliche tropische Bewegung .\ .\ . & $981''8459$ \\
\shortstack[l]{Sonnenferne; für die Epoche siderisch\\ruhend} .\ .\ .\ .\ .\ .\ .\ .\ .\ . & $69^\circ\;50'\;31''9$ \\
Aufsteigender Knoten .\ .\ .\ .\ . & $103^\circ\;18'\;28''$ \\
Neigung der Bahn .\ .\ .\ .\ .\ .\ . & $7\;8\;10.7$ \\
Excentricität .\ .\ .\ .\ .\ .\ .\ .\ .\ . & $0.0855050$ \\
Logarithm des mittlern Abstandes .\ . & $0.3720160$ \\
\end{tabular}\end{center}
\par\indent Die Differenz dieser Elemente von den zweiten ist nur unbedeutend, doch wird durch diese der
geocentrische Ort der \emph{Vesta} zur Zeit der nächsten Opposition im September 1808 um $22'$ grösser. Zugleich hatte Dr. \textsc{Gauss} die Güte, uns jene Reihe Pariser Beobachtungen nebst den Differenzen der beobachteten Positionen mit den nach den II. Elementen berechneten zu überschicken.

\smallskip
\clearpage\begingroup\fontsize{9}{11}\selectfont\setlength{\tabcolsep}{2.5pt}\renewcommand{\arraystretch}{1.05}
\begin{center}\begin{tabular}{lrrrrr}
1807 & \shortstack{Mittlere Zeit\\in Paris} & AR. \VestaRATwoNinetySix & \shortstack{Differ. d.\\II. Elem.} & Declin. \VestaDecTwoNinetySix & \shortstack{Differ. d.\\II. Elem.} \\ \hline
April 13 & $22^{\text{u}}\;39^{\text{m}}\;6\mathord{\raisebox{.65ex}{\ObservedTimeStroke}}2$ & $180^\circ 57' 7''8$ & $+3''4$ & $12^\circ 50' 35''0$ & $+5''2$ \\
14 & $22\;34\;28.3$ & $180\;46\;27.9$ & $+1.3$ & $12\;52\;39.2$ & $+6.0$ \\
18 & $22\;16\;7.0$ & $180\;\;7\;\;6.9$ & $-3.1$ & $12\;58\;\;9.3$ & $+12.7$ \\
19 & $22\;11\;34.9$ & $179\;58\;\;4.1$ & $+2.2$ & $12\;59\;\;5.5$ & $+0.4$ \\
21 & $22\;2\;39.0$ & $179\;41\;20.6$ & $-1.4$ & $12\;58\;40.7$ & $+65.0$ \\
24 & $21\;49\;20\;2$ & $179\;19\;\;5.1$ & $+1.7$ & $12\;58\;44.2$ & $+1.7$ \\
25 & $21\;44\;58.3$ & $179\;12\;26.6$ & $+4.7$ & $12\;57\;47.3$ & $+7.4$ \\
26 & $21\;40\;38.5$ & $179\;\;6\;33.8$ & $-13.1$ & $12\;56\;40.6$ & $+7.6$ \\
27 & $21\;36\;18.4$ & $179\;\;0\;35.4$ & $+0.1$ & $12\;55\;10.8$ & $+15.9$ \\
28 & $21\;32\;1.2$ & $178\;55\;15.2$ & $-0.$ & $12\;53\;47.3$ & $+2.8$ \\
29 & $21\;27\;45.5$ & $178\;50\;16.5$ & $+4.1$ & $12\;51\;54.4$ & $+4.3$ \\
30 & $21\;23\;31.7$ & $178\;45\;52.5$ & $-0.8$ & $12\;49\;49.4$ & $+3.2$ \\
Mai 1 & $21\;19\;20.0$ & $178\;41\;47.7$ & $+1.6$ & $12\;47\;26.9$ & $+5.0$ \\
7 & $20\;54\;44.7$ & $178\;26\;43.7$ & $+4.3$ & $12\;28\;39.6$ & $-7.5$ \\
12 & $20\;35\;3.3$ & $178\;26\;18.3$ & $+8.1$ & $12\;\;6\;35.8$ & $+2.8$ \\
15 & $20\;23\;35.2$ & $178\;31\;18.5$ & $+7.9$ & $11\;50\;39.4$ & $+25.5$ \\
16 & $20\;19\;50.7$ & $178\;33\;47.4$ & $+9.8$ & $11\;45\;30.8$ & $+0.2$ \\
17 & $20\;16\;5.0$ & $178\;36\;42.3$ & $+10.9$ & $11\;39\;46.8$ & $-0.6$ \\
18 & $20\;13\;22.5$ & $178\;39\;57.3$ & $+16.6$ & $11\;33\;50.4$ & $+0.3$ \\
20 & $20\;5\;1.9$ & $178\;47\;55.5$ & $+13.0$ & $11\;21\;31.8$ & $-3.5$ \\
21 & $20\;1\;24.6$ & $178\;52\;29.4$ & $+12.4$ & $11\;15\;\;5.5$ & $-3.6$ \\
22 & $19\;57\;48.9$ & $178\;57\;25.8$ & $+12.9$ & $11\;\;8\;27.6$ & $-2.1$ \\
\end{tabular}\end{center}\endgroup


\par\indent Die Zeiten sind hier von Mitternacht gezählt, und die Declination vom 21. April ist offenbar um
eine Minute zu klein. Allein wer erstaunt nicht über diese ganz vortreffliche Übereinstimmung der II.
Elemente mit dem Himmel?

\par\indent „Die Ephemeride für den Lauf der \emph{Vesta} im nächsten Jahr, schreibt uns Dr. \textsc{Gauss} ferner, folgt
hiebei. Ein grosser Theil des diesmaligen Laufs der \emph{Vesta} fällt auf die \textsc{Harding}'sche Karte für die Bewegung der \emph{Ceres} 1804, die der M. Corr. beigefügt ist. Herr Prof. \textsc{Harding} wird dasjenige Blatt seines
Atlasses, welches das übrige enthält, nächstens dem Kupferstecher übergeben. Diese nützliche Unternehmung kann nicht genug empfohlen und unterstützt werden. Das erste sehr nett gestochene Blatt
habe ich bereits in \emph{Bremen} gesehen.“

\SourceLeaf{297}\par\indent Leider wird nach beifolgender Ephemeride die \emph{Vesta} erst im Mai oder Juni des folgenden Jahres
wieder sichtbar werden. [Lauf der Vesta von 1808 Mai 2. bis 1809 März 18. mit viertägigen Zwischenräumen nach den III. Elementen AR. und Decl. auf Minuten log. des Abstandes in vier Decimalen.]

\par\indent Bei den jetzt schon ziemlich genäherten Elementen der \emph{Vesta}-Bahn wird es interessant sein, ihre
Aufsuchung in ältern Stern-Verzeichnissen zu versuchen. Unsere Bemühung eine Beobachtung der \emph{Vesta}
unter \textsc{Mayer}'s Beobachtungen aufzufinden war bis jetzt vergeblich.

\begin{flushright}
[\textsc{von Zach}.]
\end{flushright}

\Needspace{24\baselineskip}\fancyhead[L]{\small COMET 1807.}
\centerline{COMET 1807.}

\begin{center}
\rule{0.5\linewidth}{0.4pt}\\[2pt]
Monatliche Correspondenz. Band XVI. S. 562–567. 1807 December.\\[2pt]
\rule{0.5\linewidth}{0.4pt}
\end{center}
\par\indent ---\ ---\ --- Dr. \textsc{Gauss} der die Bahn keines am Himmel erscheinenden Fremdlings unbeachtet lässt,
schrieb uns unterm 16. November: „Aus Bremer und Lilienthaler Beobachtungen bis zum 5. November
habe ich folgende parabolische Elemente des Cometen abgeleitet

\smallskip

\begin{center}\begin{tabular}{lr}
Durchgang duchs Perihelium 18. Sept. .\ .\ .\ .\ . & $19^{\text{u}}\;6^{\text{m}}$ in Paris \\
Longit. Perihelii .\ .\ .\ .\ .\ .\ .\ .\ .\ .\ . & $271^\circ\;0'\;13''3$ \\
$\AscendingNode$ .\ .\ .\ .\ .\ .\ .\ .\ .\ .\ .\ .\ .\ .\ .\ . & $266\;38\;31.2$ \\
\multicolumn{2}{l}{beide vom scheinbaren Aequinoct ab gezählt} \\
Log. des kleinsten Abstandes .\ .\ .\ .\ .\ . & $9.8114927$ \\
{}[Neigung der Bahn .\ .\ .\ .\ .\ .\ .\ .\ .\ . & $63^\circ\;12'\;35''6$] \\
\multicolumn{2}{l}{Bewegung rechtläufig}
\end{tabular}\end{center}


\par\indent Diese Elemente sind nur unbedeutend verschieden von denen, die ich schon früher aus einigen
Berliner Beobachtungen berechnet hatte, und ebenso weichen sie nur wenig von denen ab, welche Hr.
\textsc{Bessel} aus einer kürzern Reihe von Beobachtungen abgeleitet hat. Von Göttingen aus werde ich Ihnen
mehr darüber schreiben.“

\begin{flushright}
[\textsc{von Zach}.]
\end{flushright}

\Needspace{15\baselineskip}\fancyhead[L]{\small JUNO. ELEMENTE VII.}
\centerline{JUNO. ELEMENTE VII.}

\begin{center}
\rule{0.5\linewidth}{0.4pt}\\[2pt]
Göttingische gelehrte Anzeigen. Stück 14. S. 129. 130. 1808 Januar 23.\\[2pt]
\rule{0.5\linewidth}{0.4pt}
\end{center}

\par\indent Von der \emph{Juno} sind aus dem verwichenen Jahre keine Beobachtungen bekannt geworden als diejenigen, welche Hr. Inspector \textsc{Bessel} in Lilienthal im März, April und Mai mit dem Kreis-Micrometer
angestellt hat. Es ist auch wahrscheinlich, dass die überaus geringe Lichtstärke dieses Planeten nirgends Beobachtungen an fixen Instrumenten erlaubt haben wird: überdies waren die meisten Astronomen zur Zeit seiner Sichtbarkeit zu sehr mit der neu entdeckten \emph{Vesta} beschäftigt, worüber die Juno
etwas vernachlässigt sein mag. Unser Hr. Prof. \textsc{Gauss} hat daher keinen Anstand genommen, jene Beobachtungen zur Verbesserung der letzten Elemente der Juno anzuwenden, zumal da Hr. \textsc{Bessel} mit
\SourceLeaf{298}einem vortrefflichen Telescop beobachtet hat, und in Behandlung des Kreis-Micrometers eine
ausgezeichnete Fertigkeit besitzt. Zuerst wurde vermittelst der Beobachtungen aus dem Mai die
Opposition bestimmt, wofür sich folgendes Resultat ergab:

\[
\begin{array}{l}
1807.\ \text{Mai}\ 17.\quad 12^{\text{u}}\,27^{\text{m}}\,41^{\text{s}}\ \text{m.\,Z. in Göttingen}\\[2pt]
\text{wahre Länge}\ .\ .\ .\ .\ .\ 236^\circ\ 4'\ 57''6.\\[2pt]
\text{wahre geocentrische Breite}\quad 16\;56\;0.1\ \text{nordl.}
\end{array}
\]\par\indent Aus der Verbindung dieser Opposition mit denen von 1804 und 1806, und andern Oertern von 1804,
1805 und 1806, wurden dann folgende neue Elemente (VII.) berechnet, wodurch die sämmtlichen
zum Grunde gelegten Beobachtungen möglichst genau dargestellt werden: [Sie stehen in diesem
Abdrucke auf Seite 294 als Theil des Briefes an \textsc{Bode} vom 24. Januar 1808.]\par\indent Die Ephemeride für die nächste Erscheinung der Juno nach diesen neuen Elementen ist bereits
vollendet, und die Astronomen werden dieselbe in der Monatlichen Correspondenz des Freiherrn
\textsc{von Zach} finden. Diese Rechnung war um so nothwendiger, da die von Hrn. Prof. \textsc{Bode} im Jahrbuche
für 1810 nach den VI. Elementen construirte Ephemeride um die Zeit der nächsten Opposition die
Rectascensionen um mehr als $\frac{3}{4}$ Grad zu gross gibt, und die Aufsuchung und
Beobachtung der Juno durch ihre auch diesmal noch sehr geringe Lichtstärke sehr erschwert
werden wird.
